\documentclass[10pt,a4paper]{amsart}
\usepackage[margin=19mm]{geometry}
\usepackage{amsmath,amssymb,mathtools}
\usepackage[scr=rsfs,cal=euler]{mathalfa}
\usepackage{xfrac}
\usepackage[unicode,hidelinks,bookmarks=false]{hyperref}
\newcommand{\ie}{i.e.\ }
\newcommand{\cf}{cf.\ }
\newcommand{\resp}{respectively\ }
\newcommand{\Subsection}{\subsection}
\providecommand{\nobreakdash}{\nobreak\hbox{-}}
\setlength{\emergencystretch}{2em}
\multlinegap0pt
\usepackage{amssymb}
\newtheorem{thm}{Theorem}
\newtheorem{pf}{Pf}
\newcommand{\norm}[1]{\left\lVert#1\right\rVert}
\title{Multidimensional Bohr radii for holomorphic functions with values in complex Banach spaces}
\author{Vasudevarao Allu, Himadri Halder, Subhadip Pal}
\date{Source: arXiv:2308.07825v5 (CC BY 4.0)}
\begin{document}
\maketitle

\noindent\emph{Source and licence.} Multidimensional Bohr radii for holomorphic functions with values in complex Banach spaces. Version arXiv:2308.07825v5 (CC BY 4.0), distributed by arXiv under the Creative Commons Attribution 4.0 International licence, http://creativecommons.org/licenses/by/4.0/, as recorded in the arXiv licence metadata for this exact version and shown on the abstract page https://arxiv.org/abs/2308.07825v5. Full licence notice: \texttt{licenses/arxiv-2308.07825-CC-BY-4.0.txt}.\par\smallskip

\noindent\textit{Editorial scope.} Counted unit: the article's thm Pal-Vasu-P3-thm-2.6 (source0.tex lines 424--429). The article's complete original proof of it is delivered with it (source0.tex lines 610--654, one \texttt{proof} environment of 45 lines, which the article places away from the statement). No mathematics is written, repaired or re-derived: the statement and the proof are the article's own lines, in the article's own notation, and no retained line is substituted or re-flowed.  No further source-local context is needed: every cross-reference of every retained line resolves inside the two delivered spans.  Nothing was omitted from any retained span, so the package carries no omission record.  No retained line contains a citation, so the wrapper carries no bibliography and no bibliography entry.  No retained line contains a figure environment, an image-inclusion command or drawing code, so no image is delivered and the package declares no asset dependency.  Editorial formatting only, in addition: an AMS \texttt{amsart} preamble that restates the article's own declarations and macros used by the retained lines, namely the environments \texttt{thm}, \texttt{pf} on separate counters, together with the standard packages those lines need.  The retained environments print in this document as Theorem 1, Pf 1 in order of appearance, whereas the article prints the same items with the numbering of its own full document.  The complete native source of the article as distributed in the arXiv e-print for this exact version is bundled unchanged as \texttt{sources/145679/source0.tex}.
\begin{thm}\label{Pal-Vasu-P3-thm-2.6}
	Let $\Omega$ be a complete Reinhardt domain in $\mathbb{C}^n$ and $\lambda\geq 1$. If $\mathcal{O}(\Omega, X)$ is a space of all holomorphic mappings on $\Omega$ into a complex Banach space $X$ containing all vector-valued polynomials then for every $1\leq p,q, <\infty$, we have
	\begin{equation*}
		A_{p,q}(\mathcal{O}(\Omega, X), \lambda)=A_{p,q}(H_{\infty}(\Omega, X), \lambda).
	\end{equation*}
\end{thm}

\begin{pf} [{\bf Proof of Theorem \ref{Pal-Vasu-P3-thm-2.6}}]
	We start with an observation that 
	\begin{equation*}
		\mathcal{P}(\Omega, X)\subseteq \mathcal{O}(\Omega, X)\subseteq H_{\infty}(\Omega, X).
	\end{equation*}
	Then, we have the following inclusion
	\begin{equation*}
		A_{p,q}(H_{\infty}(\Omega, X), \lambda) \leq A_{p,q}(\mathcal{O}(\Omega, X), \lambda) \leq A_{p,q}(\mathcal{P}(\Omega, X), \lambda).
	\end{equation*}
	It suffices to show that for any $\lambda\geq 1$, $A_{p,q}(\mathcal{P}(\Omega, X), \lambda)\leq 	A_{p,q}(H_{\infty}(\Omega, X), \lambda)$.\\[2mm]
	Assume $r\in \mathbb{R}^n_{\geq 0}$ such that for all $P\in \mathcal{P}(\mathbb{C}^n, X)$,
	\begin{equation}\label{Pal-Vasu-P3-e-2.14}
		\norm{x_0(P)}^p+\left(\sum_{k=1}^{\infty}\norm{x_{\alpha}(P)}r^{\alpha}\right)^q\leq \lambda \norm{P}_{\Omega}.
	\end{equation}
	Fix $f\in H_{\infty}(\Omega, X)$ and $0<t<1$. Our aim is to show that 
	\begin{equation}\label{Pal-Vasu-P3-e-2.15}
		\norm{x_0(f)}^p+\left(\sum_{k=1}^{\infty}\norm{x_{\alpha}(f)}(tr)^{\alpha}\right)^q\leq \lambda \norm{f}_{\Omega}.
	\end{equation}
	For $k\geq 0$, consider the $k$-homogeneous polynomials $\chi_{k}=\sum_{|\alpha|=k}x_{\alpha}(f)z^{\alpha}$, where $x_{\alpha}(f)$ is the $\alpha$-th coefficient in the monomial series expansion $\sum_{\alpha\in \mathbb{N}^n_{0}}x_{\alpha}(f)z^{\alpha}$ of $f$. Note that $\Omega$ is balanced, so by Cauchy inequalities we have 
	\begin{equation}\label{Pal-Vasu-P3-e-2.16}
		\norm{\chi_{k}}_{t\Omega}=t^k\norm{\chi_{k}}_{\Omega}\leq t^k \norm{f}_{\Omega}.
	\end{equation}
	Therefore, the series $\sum_{k=0}^{\infty}\chi_{k}$ converges absolutely and uniformly to $f$ on $t\Omega$. Thus, given $\epsilon>0$, there exits $n_0 \in \mathbb{N}$ such that for all $n \geq n_0$,
	\begin{equation*}
		\bigg\|\sum_{k=0}^{n}\chi_{k}\bigg\|_{t\Omega}\leq \norm{f}_{t\Omega} + \epsilon \leq \norm{f}_{\Omega}+\epsilon.
	\end{equation*}
	We now consider the polynomials $\Gamma_{k}(z)=\chi_{k}(tz)$ for $z\in \mathbb{C}^n$. Then,
	\begin{equation*}
		x_{\alpha}(\Gamma_{k})=
		\left\{\begin{array}{ll}
			t^kx_{\alpha}(f) & \mbox{ if  \,$|\alpha|=k$},\\[1mm]
			0  & \mbox{ if \,$|\alpha|\neq k$}.
		\end{array}\right.
	\end{equation*}
	From \eqref{Pal-Vasu-P3-e-2.14} and \eqref{Pal-Vasu-P3-e-2.16}, we obtain that
	\begin{align*}
		\norm{x_0(f)}^p+\left(\sum_{k=1}^{\infty}\norm{x_{\alpha}(f)}(tr)^{\alpha}\right)^q &= \norm{x_0(\Gamma_{0})}^p + \left(\sum_{\alpha\in \mathbb{N}^n}\bigg\|x_{\alpha}\left(\sum_{k=1}^{n}\Gamma_{k}\right)r^{\alpha}\bigg\|\right)^q\\[2mm] & \leq \lambda \bigg\|\sum_{k=0}^{n}\Gamma_{k}\bigg\|_{\Omega}=\lambda \bigg\|\sum_{k=0}^{n}\chi_{k}\bigg\|_{t\Omega}\leq \lambda \left(\norm{f}_{\Omega}+\epsilon\right),
	\end{align*} 
	which holds for all $\epsilon>0$. Consequently, this proves \eqref{Pal-Vasu-P3-e-2.15}.
	Hence for every $r\in \mathbb{R}^n_{\geq 0}$ satisfying \eqref{Pal-Vasu-P3-e-2.14} and for all $f\in H_{\infty}(\Omega, X)$ and $0<t<1$, we have shown that 
	\begin{equation*}
		t\sum_{i=1}^{n}r_i \leq A_{p,q}(H_{\infty}(\Omega,X),\lambda),
	\end{equation*}
	which proves our claim. This completes the proof.
\end{pf}

\end{document}
